\chapter{Maxillary Le Fort 1 osteotomy}

\section*{Definition and Overview}
A Maxillary Le Fort 1 osteotomy is a major reconstructive craniofacial and orthognathic surgical procedure characterised by a horizontal fracture-line osteotomy across the maxilla, separating the tooth-bearing alveolar segment of the upper jaw from the upper craniofacial skeleton. This procedure allows the surgeon to reposition the maxilla in three dimensions (advancement, intrusion, extrusion, setback, or rotation) to correct dentofacial deformities, malocclusion, and obstructive sleep apnoea. It is performed entirely intraorally, eliminating external facial scars. Repositioning the maxilla restores facial harmony, improves masticatory function, enhances airway patency, and establishes an aesthetic occlusion in conjunction with orthodontic therapy.

\section*{Epidemiology}
Orthognathic surgery, including the Le Fort 1 osteotomy, is commonly performed in adolescents and young adults once skeletal maturity is complete (usually between 16 to 18 years in females and 18 to 21 years in males). In Australia, thousands of orthognathic procedures are conducted annually by oral and maxillofacial surgeons in private and public hospital settings. It is indicated for patients with moderate-to-severe skeletal discrepancies (Class II or Class III malocclusions, vertical maxillary excess, or skeletal open bites) that cannot be resolved with orthodontics alone.

\section*{Aetiology and Pathophysiology}
The conditions necessitating a Le Fort 1 osteotomy are primarily structural skeletal discrepancies rather than active infectious or inflammatory processes.
\begin{itemize}
    \item \textbf{Skeletal discrepancies:} Genetic or developmental imbalances in jaw growth, such as maxillary hypoplasia (underdevelopment of the upper jaw, common in Class III malocclusion) or vertical maxillary excess (excessive vertical growth, leading to a ``gummy smile'').
    \item \textbf{Cleft lip and palate:} Congenital craniofacial anomalies often result in secondary maxillary hypoplasia and scarring, necessitating osteotomy to advance the maxilla.
    \item \textbf{Obstructive sleep apnoea (OSA):} Skeletal retrognathia narrows the posterior airway space; advancing the maxilla (often combined with mandibular advancement) widens the airway and reduces pharyngeal collapse.
    \item \textbf{Surgical anatomy:} The osteotomy line runs horizontally above the dental roots, through the nasal septum, lateral nasal walls, and pterygoid plates. This requires preservation of the descending palatine arteries to maintain blood supply to the mobilised maxillary segment.
\end{itemize}

\section*{Clinical Features}
Patients requiring a Le Fort 1 osteotomy present with distinct aesthetic and functional signs.
\begin{itemize}
    \item \textbf{Malocclusion:} Inability of the upper and lower teeth to meet properly (e.g.\ anterior open bite, crossbite, or severe overjet).
    \item \textbf{Masticatory dysfunction:} Difficulty chewing or biting certain foods, often leading to temporomandibular joint (TMJ) discomfort or strain.
    \item \textbf{Facial disharmony:} Features such as a retruded midface, flat cheekbones, or an excessively long face with a receding chin.
    \item \textbf{Airway restriction:} Snoring, mouth breathing, daytime sleepiness, and documented obstructive sleep apnoea due to pharyngeal airway narrowing.
    \item \textbf{Speech articulation issues:} Difficulties pronouncing certain sounds (lisping) due to abnormal jaw and dental relationships.
\end{itemize}

\section*{Differential Diagnosis}
Skeletal maxillomandibular discrepancies must be differentiated from dental-only malocclusions and other craniofacial conditions.
\begin{itemize}
    \item \textbf{Dental malocclusion:} Teeth are misaligned but the underlying skeletal jaw relationship is normal; this can be treated with orthodontics alone without surgical intervention.
    \item \textbf{Le Fort 2 or 3 osteotomy:} Indicated for severe midface hypoplasia involving the nasal bones (Le Fort 2) or the zygomatic arches and orbits (Le Fort 3), commonly seen in syndromic craniosynostosis (e.g.\ Crouzon or Apert syndrome).
    \item \textbf{Distraction osteogenesis:} A technique where the bone is cut and slowly distracted (pulled apart) over time to grow new bone, preferred in patients with severe tissue deficiency or in younger, growing patients.
    \item \textbf{Isolated mandibular surgery:} Repositioning only the lower jaw (sagittal split osteotomy) when the maxilla is correctly positioned but the mandible is hypoplastic or hyperplasia-dominant.
\end{itemize}

\section*{When to Seek Medical Care}
Orthognathic surgery is an elective, planned procedure. However, patients recovering from a Le Fort 1 osteotomy must monitor for post-operative complications.

\textbf{Call 000 (or local emergency services) for:}
\begin{itemize}
    \item Severe, uncontrollable bleeding from the mouth or nose that does not stop with light pressure.
    \item Acute breathing difficulties, airway obstruction, or sudden swelling of the throat or tongue.
    \item Sudden collapse, severe dizziness, or loss of consciousness.
\end{itemize}
Other reasons to see a doctor: Post-operative patients should contact their surgeon if they develop a high fever, worsening pain, an unstable bite, or signs of wound infection (such as foul-smelling discharge or spreading redness inside the mouth).

\section*{Diagnosis and Investigations}
Pre-operative planning is highly detailed, combining clinical assessment with three-dimensional imaging and computer-aided design.
\begin{itemize}
    \item \textbf{Cone-beam computed tomography (CBCT):} Provides three-dimensional skeletal views of the craniofacial complex, allowing precise measurement of bone density, volume, and airway space.
    \item \textbf{Digital orthodontic models:} Intraoral scans of the teeth used to plan the post-surgical occlusion and fabricate surgical splints.
    \item \textbf{Virtual surgical planning (VSP):} Computer software simulating the osteotomy, movement of the maxillary segment, and pre-fabrication of custom titanium plates and screws.
    \item \textbf{Lateral cephalometric analysis:} Standardised skull X-rays used to measure jaw angles, dental inclinations, and relationship to the cranial base.
    \item \textbf{Nasendoscopy and sleep studies:} In patients with sleep apnoea, to evaluate the site of airway obstruction and establish a baseline apnoea-hypopnea index (AHI).
\end{itemize}

\section*{Management}
Management involves a coordinated, multidisciplinary approach spanning pre-operative preparation, the surgical procedure, and post-operative rehabilitation.
\begin{itemize}
    \item \textbf{Combined orthodontic-surgical treatment:} Typically involves 12 to 18 months of pre-surgical orthodontics to align the teeth within each arch, followed by post-surgical orthodontics to fine-tune the bite.
    \item \textbf{Surgical procedure (Le Fort 1):} Performed under general anaesthesia with nasal intubation. An incision is made in the upper gum line, the bone is cut, repositioned using a prefabricated surgical splint, and rigidly fixed with titanium plates and mini-screws.
    \item \textbf{Post-operative dietary modifications:} A strict liquid-to-soft diet for 6 weeks to allow the osteotomy sites to heal without skeletal displacement.
    \item \textbf{Pain and swelling management:} Use of intravenous and oral analgesics, cold compresses to the cheeks, head elevation, and occasionally corticosteroids to minimise facial swelling.
    \item \textbf{Elastics and jaw tracking:} Use of light orthodontic elastics to guide the jaw into the new bite, avoiding rigid intermaxillary fixation (wiring the jaws shut) where possible.
\end{itemize}

\section*{Complications}
While Le Fort 1 osteotomy is highly successful, it carries inherent risks associated with major craniofacial surgery.
\begin{itemize}
    \item \textbf{Post-operative haemorrhage:} Bleeding from the maxillary artery or its branches, which may require nasal packing or, rarely, embolisation.
    \item \textbf{Sensory nerve deficit:} Numbness or paraesthesia of the infraorbital nerve distribution (cheek, upper lip, gums, and teeth), which is usually temporary but can be permanent in a small percentage of cases.
    \item \textbf{Infection or hardware failure:} Infection around the titanium plates or loosening of screws, requiring a course of antibiotics or removal of hardware.
    \item \textbf{Avascular necrosis of the maxilla:} A rare, severe complication where blood supply to the mobilised maxilla is compromised, leading to bone or tooth loss.
    \item \textbf{Relapse:} Gradual shifting of the jaw back toward its original position, necessitating secondary orthodontic or surgical correction.
    \item \textbf{Nasal changes:} Widening of the nasal base (alar flare) or deviation of the nasal septum, which surgeons attempt to control using an alar cinch suture and septal shaping.
\end{itemize}

\section*{Prognosis and Outcomes}
The prognosis for patients undergoing a Maxillary Le Fort 1 osteotomy is highly positive, with high rates of patient satisfaction regarding both facial aesthetics and functional outcomes. Bone healing (union) is typically complete within 6 to 8 weeks, allowing a return to a normal diet. Orthodontic treatment is completed within 6 months post-surgery. Long-term stability of the repositioned maxilla is excellent, particularly when rigid internal fixation and appropriate post-surgical retention are utilised.

\section*{Prevention and Self-Care}
Post-operative self-care is critical to preventing infection and ensuring proper skeletal healing.
\begin{itemize}
    \item \textbf{Maintain strict oral hygiene:} Rinse the mouth gently with warm salt water or chlorhexidine mouthwash after every meal, and use a soft-bristled baby toothbrush to clean the teeth.
    \item \textbf{Adhere to diet restrictions:} Strictly follow the soft diet guidelines (no chewing) for the recommended 6 weeks.
    \item \textbf{Avoid nose blowing:} Do not blow the nose or sneeze with the mouth closed for 4 weeks post-surgery to prevent air forcing into the sinuses and facial tissues.
    \item \textbf{Elevate the head:} Sleep with the head elevated on 2 to 3 pillows for the first 2 weeks to reduce facial swelling.
    \item \textbf{Limit physical exertion:} Avoid strenuous activities, heavy lifting, and contact sports for at least 8 to 12 weeks to protect the healing facial bones.
\end{itemize}

\section*{Sources and Further Reading}
The following professional bodies and resources provide reliable information on orthognathic surgery:
\begin{itemize}
    \item Australian Society of Orthodontists. \textit{Patient guides on combined orthodontic and jaw surgery treatments.} https://www.aso.org.au/
    \item Royal Australasian College of Dental Surgeons. \textit{Information on training and clinical standards in oral and maxillofacial surgery.} https://www.racds.org/
    \item British Association of Oral and Maxillofacial Surgeons. \textit{Patient information leaflets on Le Fort 1 maxillary osteotomy.} https://www.baoms.org.uk/
    \item American Association of Oral and Maxillofacial Surgeons. \textit{Resources on orthognathic surgery procedures and patient recovery.} https://www.myoms.org/
    \item Mayo Clinic. \textit{Overview of orthognathic surgery, including what to expect during and after jaw surgery.} https://www.mayoclinic.org/tests-procedures/jaw-surgery/about/pac-20384990
\end{itemize}
